\documentclass[11pt,a4paper]{article}

% Companion draft to paper-vector-navigation.tex; venue-neutral.
\usepackage[T1]{fontenc}
\usepackage[utf8]{inputenc}
\usepackage{amsmath,amssymb,amsthm}
\usepackage{xcolor}
\usepackage[colorlinks=true,linkcolor=blue!50!black,citecolor=blue!50!black,urlcolor=blue!50!black]{hyperref}
\usepackage{microtype}
\usepackage{fullpage}
\usepackage{setspace}
\onehalfspacing
\emergencystretch=1.5em
\usepackage{enumitem}
\setlist{itemsep=2pt,topsep=4pt}

\newtheorem{definition}{Definition}[section]
\newtheorem{property}{Property}[section]
\newtheorem{guarantee}{Guarantee}[section]
\theoremstyle{remark}
\newtheorem{remark}{Remark}[section]

\newcommand{\code}[1]{\texttt{#1}}

\title{The Client Is the Database:\\ Caching, Consistency, and Projection Maintenance\\ for Applications over Personal Data Stores}
\author{Draft --- author list TBD}
\date{2026-07-12 (working draft)}

\begin{document}
\maketitle

\begin{abstract}
In an application over personal data stores (Solid pods), the storage layer is a
set of plain HTTP resource servers the application does not control: there is no
server-side query engine, no reliable push channel, and several agents write. The
layers a database management system normally provides --- concurrency control,
cache coherence, materialised-view maintenance --- must therefore be
reconstructed client-side over HTTP. We describe such a reconstruction. Writes
commit against one of two \emph{storage models} --- in-place resources under
optimistic concurrency, or append-only event logs whose current state is a
\emph{projection} --- with the access-control lists maintained as a replayable
materialised projection of a user-owned grant log. Reads flow through a
\emph{resource-oriented} client cache that mirrors the store's container/resource
graph, keyed by IRI, deriving screen shapes at the edge; we show this keying makes
cache coverage structural (a derived read cannot silently under-cover its inputs).
Freshness is maintained without push by four mechanisms --- server-driven
revalidation, write-owned invalidation, last-good-wins rendering, and
\emph{attention-triggered} revalidation for state written by other agents --- and
we state the resulting consistency contract in the style of session guarantees.
Finally we report conformance findings from running against three Solid server
implementations: entity validators are missing or, in one case, unsound under
content negotiation, and retry layers must sit above per-request proof-of-possession
signing. The system is in production use in a building-energy application; a
companion paper describes the interaction model above this layer.
\end{abstract}

\section{Introduction}\label{sec:intro}

The premise of personal data stores is that applications \emph{visit} data that
users own~\cite{solid}. Architecturally this deletes the database tier: a Solid
pod is a plain HTTP resource server speaking LDP~\cite{ldp} --- \code{GET},
\code{PUT}, \code{POST}, \code{DELETE} over RDF documents and containers --- with
per-resource access control~\cite{wac}. Three consequences follow for any
non-trivial application:

\begin{itemize}
\item \textbf{No query engine.} The store answers only dereference and container
  listing. Joins, folds, and projections happen in the client.
\item \textbf{No push.} Change notification exists as a
  specification~\cite{solid-notifications} but cannot be assumed across servers;
  in practice the client polls, and must decide \emph{when}.
\item \textbf{Multiple writers.} The user's own devices, the application's own
  reconciliation, and --- through sharing --- \emph{other agents} write state the
  client renders. Some of that state (another user's access-control lists) the
  client cannot even read.
\end{itemize}

A DBMS would meet these with concurrency control, cache coherence, and
materialised-view maintenance. Our thesis: \textbf{the client must become its own
database}, and doing so systematically --- rather than screen by screen --- is
what keeps a decentralised application correct. This paper reports the design as
built and hardened in Granergize, a production React application for
building-energy data on Solid pods, and makes four contributions:

\begin{enumerate}
\item \textbf{A write-side taxonomy} (\S\ref{sec:writes}): two storage models
  (in-place vs event-sourced) chosen by where authority lives, two projection
  disciplines (fold-on-read vs materialised) chosen by whether the consumer can
  fold, and an orthogonal trigger axis (user-intent vs reconciliation) --- with
  the WAC access-control lists as a \emph{replayable materialised projection} of
  a user-owned grant log, including the audit and repair operations that keep it
  honest, and an idempotency discipline (client-chosen URIs, conditional
  \code{PUT}) that makes creation replay-safe under a retrying transport.
\item \textbf{A resource-oriented client cache} (\S\ref{sec:reads}): entries
  mirror the store's container/resource graph, keyed by IRI and namespaced by
  viewer; screen shapes are derived at the edge by pure selectors. We state the
  structural property this buys: a derived read cannot under-cover its inputs,
  because the unit of invalidation is the resource a write touches.
\item \textbf{A consistency contract without push} (\S\ref{sec:consistency}):
  four freshness mechanisms, including \emph{attention-triggered revalidation}
  for cross-agent state, and the resulting guarantees stated in the style of
  session guarantees~\cite{terry1994} --- read-your-writes for own-pod state,
  monotonic display per session, attention-bounded eventual consistency across
  agents, and self-healing visibility under revocation.
\item \textbf{Ecosystem conformance findings} (\S\ref{sec:findings}): what three
  Solid server implementations actually provide of the HTTP machinery this
  design leans on --- missing entity validators, a content-negotiation
  unsoundness in conditional requests, and the interaction of retries with
  DPoP-style proof-of-possession --- plus the in-application instrumentation used
  to find them.
\end{enumerate}

A companion paper describes the interaction model \emph{above} this layer (a
typed command--query intent catalogue and a navigation state machine); the
interface between the two is exactly the intent dispatch: this paper is what
happens beneath a read intent's answer and a write intent's commit.

\section{A running example}\label{sec:running}

Alice operates logistics halls; her pod holds one RDF document per building and
one per energy year. She shares the Feucht hall --- energy included, years
2022--2024 --- with her consultant Carol. Three things must now stay true, and
each is threatened by a different failure mode:

\begin{enumerate}
\item \emph{Carol can read exactly what Alice granted.} The grant's scope lives
  in an event Alice's pod records; the server enforces only \code{.acl} files.
  If the two drift --- an archive restore that recovers the log but not the
  ACLs, or a new 2025 energy year created \emph{inside} an all-years grant ---
  enforcement silently diverges from intent (\S\ref{sec:projections}).
\item \emph{Alice's screens reflect her writes, immediately.} She edits the hall
  and saves a reading; every view folding that document must refetch --- and
  nothing else should (\S\ref{sec:reads}).
\item \emph{Alice eventually sees what others did.} Bob revokes a share to her;
  a data-room member joins. No notification may arrive; her client must decide
  when to look (\S\ref{sec:consistency}).
\end{enumerate}

The design below is organised around exactly these three obligations.

\section{Setting and requirements}\label{sec:setting}

The store is a set of pods: per-user HTTP servers holding RDF (Turtle) documents
in LDP containers, authenticated by Solid-OIDC with DPoP proof-of-possession, and
authorised per resource by WAC \code{.acl} documents that only the server's
enforcement engine reads. The application reads three \emph{tiers} of data ---
\emph{owned} (the user's pod), \emph{granted} (peers' pods, via shares), and
\emph{open} (public registers, spatially gated) --- and the reachable data set is
assembled at read time; there is no global index, so no read may claim
exhaustiveness. From the failure modes of \S\ref{sec:running} we take the
requirements: (R1) writes must be safe under concurrency \emph{and} under a
retrying transport; (R2) enforcement state must remain rebuildable from recorded
intent; (R3) a write must refresh exactly the reads it affects; (R4) state
written by other agents must become visible without a push channel; (R5) all of
it must degrade gracefully on servers that implement the HTTP machinery
partially (\S\ref{sec:findings}).

\paragraph{The store, concretely.} All application state lives under one
\emph{app collection} on the user's storage, \code{granergize/} under the
pod's storage root; Figure~\ref{fig:podlayout} shows its layout. Container and resource URIs are
client-chosen (\S\ref{sec:idempotency}) and structure-bearing: a building's
attachments and its energy datasets are derivable from the building's own URI
without a discovery read, and the observation tree is \emph{time-first}
(year before entity) so a period's data is one container listing. This graph
is exactly what the client cache of \S\ref{sec:reads} mirrors, entry per
resource. Table~\ref{tab:statekinds} classifies each resource kind along the
dimensions \S\ref{sec:writes} develops --- which storage model governs it, who
writes it, and how its current state is derived.

\begin{figure}[t]
\begin{center}
\small
\begin{tabular}{@{}l@{\qquad}l@{}}
\texttt{granergize/} & \\
\texttt{|--~buildings/} & \\
\texttt{|~~~|--~\{id\}.ttl} & building master data (subject \texttt{<\#it>}) \\
\texttt{|~~~`--~\{id\}/files/\{name\}} & binary attachments \\
\texttt{|--~observations/} & \\
\texttt{|~~~|--~\{year\}/\{id\}.ttl} & annual energy datasets (time-first) \\
\texttt{|~~~`--~\{year\}/\{month\}/\{day\}/\{id\}.ttl} & sub-daily series \\
\texttt{|--~aggregations/} & \\
\texttt{|~~~|--~\{id\}.ttl} & roll-up definitions \\
\texttt{|~~~`--~snapshots/\{id\}.ttl} & computed snapshots \\
\texttt{|--~rooms/} & \\
\texttt{|~~~`--~\{id\}/\ldots} & data-room event logs (shared) \\
\texttt{|--~shared-out/} & outgoing grant log \\
\texttt{|--~shared-in/} & incoming grant log \\
\texttt{|--~inbox/} & notification inbox (discoverable) \\
\texttt{|--~prefs.ttl} & preferences \\
\texttt{|--~bookmarks.ttl} & room bookmarks \\
\texttt{`--~agents.ttl} & address book \\
\end{tabular}
\end{center}
\caption{The app collection on a user's storage (a relative reference,
resolved against the pod's storage root). Every entry is an ordinary LDP
resource or container; the tree is what the resource-oriented cache
(\S\ref{sec:reads}) mirrors. Room containers are the one place other agents
write inside a user's collection (via the room's own ACL); the inbox receives
\code{POST}s from peers and is drained into \code{shared-in/}. The
per-resource \code{.acl} documents sit beside the resources they protect and
are omitted here; Figure~\ref{fig:acls} gives their shape.}
\label{fig:podlayout}
\end{figure}

\begin{figure}[t]
\centering
\small
\setlength{\tabcolsep}{4pt}
\begin{tabular}{@{}lll@{}}
\code{.acl} of & authorization (agent) & modes \\
\hline
any owned resource & \code{\#owner} (owner WebID) & \code{Read},
  \code{Write}, \code{Control} \\[3pt]
shared building \code{\{id\}.ttl} & \code{\#owner}; + one per grantee WebID
  & grantee: \code{Read} \\
its \code{files/} container & \code{\#owner}; + grantee, with
  \code{acl:default} & grantee: \code{Read} \\
each in-scope energy dataset & \code{\#owner}; + grantee (per granted year)
  & grantee: \code{Read} \\[3pt]
room container \code{rooms/\{id\}/} & \code{\#owner} (creator);
  \code{\#members} & members: \code{Read}, \code{Append} \\
 & \quad(= \code{acl:AuthenticatedAgent}); both & \\
 & \quad carry \code{acl:default} & \\
inbox \code{inbox/} & \code{\#owner}; \code{\#append}
  (\code{acl:AuthenticatedAgent}) & peers: \code{Read}, \code{Append} \\
\end{tabular}
\caption{The enforcement cache: the authorizations each \code{.acl} carries.
Every owned resource grants its owner full control; a share adds one
\code{acl:Read} authorization per grantee to the building document, its
\code{files/} container (with \code{acl:default}, so later uploads inherit),
and each energy dataset the grant's recorded scope covers. Room containers
and the inbox are the two places other agents may write: members of a room
read and append under its container-wide default; any authenticated agent
may deliver into the inbox. All of these are derived state --- rebuilt from
the \code{shared-out/} log by replay (\S\ref{sec:projections}).}
\label{fig:acls}
\end{figure}

\begin{table}[t]
\centering
\small
\setlength{\tabcolsep}{3.5pt}
\begin{tabular}{@{}llll@{}}
resource kind & storage model & writer(s) & current state \\
\hline
building doc, energy dataset & in-place, cond.\ \code{PUT} & owner & the resource itself \\
attachments (binary) & in-place, immutable & owner & the resource itself \\
aggregation defs + snapshots & in-place, cond.\ \code{PUT} & owner & the resource itself \\
prefs, bookmarks, address book & in-place, cond.\ \code{PUT} & owner & the resource itself \\
outgoing share log & event log, \code{POST} & owner & fold-on-read \\
incoming share log & event log, \code{POST} & owner (drain) & fold-on-read \\
room logs & event log, \code{POST} & all participants & fold-on-read \\
inbox & transport, \code{POST} & other agents & drained, not folded \\
\code{.acl} files & materialised projection & this client (replayable) & rebuilt from grant log \\
\end{tabular}
\caption{Every state kind, classified along the write-side dimensions of
\S\ref{sec:writes}: which storage model governs it, who writes it, and how its
current state is derived. The one materialised projection is the enforcement
cache (\S\ref{sec:projections}); everything else is either ground truth
in place or a log folded on read.}
\label{tab:statekinds}
\end{table}

\paragraph{Documents and their classes.} Each resource is a Turtle document;
Table~\ref{tab:classes} lists the principal RDF classes per document kind.
The vocabulary mix is standard-first: buildings are RealEstateCore, raw
readings SOSA, aggregates the RDF Data Cube, grants the Solid interop
vocabulary, room events ActivityStreams, agents FOAF/vCard, enforcement WAC,
lineage PROV --- with three small application-owned namespaces
(\code{bldg:}, \code{cons:}, \code{gran:}) for what no standard covers
(building subsystems, the energy-dataset envelope, app preferences). The
schema is enforced by the application's parsers and serialisers, not by a
shapes language: no ShEx or SHACL artefact exists, and a document with
unexpected triples is read leniently (unknown triples survive a
read--modify--write untouched). For this paper the class structure matters
mainly as the unit of parsing: a cache entry (\S\ref{sec:reads}) holds one
document parsed against these classes. Instances follow two idioms, and the
table's middle column makes them explicit: an \emph{addressable entity} is a
fragment of its document (a building is \code{\ldots\{id\}.ttl\#it}, a
dataset \code{\ldots\{id\}.ttl\#ds}) --- the document is the unit of caching
and concurrency, the fragment the entity --- whereas an \emph{immutable
event} is its document (\code{<> a interop:AccessGrant}), one event per
server-minted log resource, which is what lets event entries opt out of
revalidation entirely (\S\ref{sec:mechanisms}).

\paragraph{The open tier, concretely.} The granted tier reuses the layout
above --- a peer's app collection, read with the recipient's own credentials
--- but the open tier is different in kind: external Linked-Data services,
resolved from a static registry of wrapper bases
(\code{https://wunderfacts.com/\{mastr, dwd, regionalstatistik, energieatlas,
lod2-by, nuts, lau, netztransparenz, addressapi\}/}, plus Wikidata and
Wikimedia Commons for operator identities and logos), reached over plain
unauthenticated HTTP through a second, read-only transport --- retried and
instrumented like the authed one, but with no DPoP and no identity, so it can
be ambient where the Pod gateway must be threaded. Discovery is query-shaped:
each wrapper exposes a small capability vocabulary over its base ---
dereference plus \code{/search?q=}, \code{/within?bbox=W,S,E,N},
\code{/nearby?lon=\&lat=\&r=}, \code{/contains?lat=\&lon=},
\code{/filter?\{attrs\}} --- and these query URLs are exactly the
\emph{spatial gates} of the tier model: the reachable open set is what the
gates issued so far have returned, never a corpus. Responses are negotiated
per source --- Turtle for records and cubes, GeoJSON for region boundaries;
Table~\ref{tab:openclasses} lists the principal classes. Open entries share
the client cache but not its freshness contract: no local write ever touches
them, so they are never write-invalidated (\S\ref{sec:mechanisms});
keyed by external IRI or gate URL, they instead carry long time-based
staleness --- a region's boundary or a published statistic changes on
bureaucratic, not interactive, timescales.

\begin{table}[t]
\centering
\small
\setlength{\tabcolsep}{4pt}
\begin{tabular}{@{}llp{6.2cm}@{}}
source & discovery & principal classes / format \\
\hline
\code{regionalstatistik/}, & filter, deref &
  \code{qb:DataSet} cubes of \code{qb:Observation}s; \code{skos:Concept}
  dimension code lists \\
\quad\code{energieatlas/} & & \\
\code{dwd/} (weather) & \code{near?lat\ldots}, \code{values?station} &
  \code{sosa:Observation} series, \code{sosa:Sensor} stations \\
\code{mastr/}; \code{netztransparenz/} & within (bbox), filter, deref &
  typed generation-unit records (\code{mastr:} vocabulary); per-plant EEG
  settlement figures \\
\code{lod2-by/} & filter, deref & \code{lod2:Building},
  \code{lod2:RoofSurface}, \code{lod2:GroundSurface} \\
\code{lau/}, \code{nuts/} & contains, search, deref & \code{skos:Concept}
  region hierarchy; boundary polygons as GeoJSON \\
\code{addressapi/} & \code{search?\{road,postcode,\ldots\}} & register
  address points (JSON) \\
\end{tabular}
\caption{The open tier: external sources (bases under
\code{https://wunderfacts.com/}), their discovery gates, and the principal
classes the client parses from them. All are read over unauthenticated HTTP
with per-source content negotiation; none is ever write-invalidated.}
\label{tab:openclasses}
\end{table}

\begin{table}[t]
\centering
\small
\setlength{\tabcolsep}{4pt}
\begin{tabular}{@{}llp{6.4cm}@{}}
document kind & instance IRI & principal classes \\
\hline
building & \code{buildings/\{id\}.ttl\#it} & \code{rec:Building}; subsystem
  fragment nodes \code{bldg:PVSystem}, \code{bldg:BatteryStorage},
  \code{bldg:CHPSystem}, heating (\code{bldg:GasBoiler},
  \code{bldg:HeatPump}, \ldots); \code{bldg:Address}, \code{geo:Point};
  \code{schema:MediaObject} (attachment links); \code{prov:Attribution} \\
energy dataset & \code{observations/\ldots/\{id\}.ttl\#ds} &
  \code{cons:EnergyDataset} (dual-typed \code{sosa:ObservationCollection}),
  one per building, year, and scenario (\code{cons:Actual} vs
  \code{cons:Planned}); members \code{sosa:Observation}; rolled-up figures
  \code{qb:Observation} \\
aggregation & \code{aggregations/\{id\}.ttl},
  & \code{cons:AggregationDefinition}; \code{cons:AggregationSnapshot},
  \code{cons:BenchmarkResult} \\
 & \code{\ldots/snapshots/\{id\}.ttl} & \\
grant event & \code{shared-out/\{event\}} $=$ \code{<>} &
  \code{interop:AccessGrant}, \code{interop:AccessRevocation} (mirrored into
  \code{shared-in/} on the recipient) \\
room event & \code{rooms/\{id\}/\{event\}} $=$ \code{<>} & \code{as:Join},
  \code{as:Leave}, \code{as:Update}; \code{gran:UserRole} \\
address book & \code{agents.ttl\#book}; contacts by WebID &
  \code{vcard:AddressBook}; \code{vcard:Individual},
  \code{vcard:Organization} (also \code{foaf:Agent} in grants) \\
singletons & \code{prefs.ttl}, \code{bookmarks.ttl} $=$ \code{<>} &
  \code{gran:Preferences}, \code{gran:Bookmarks} \\
\code{.acl} & \code{\{resource\}.acl\#owner}, \code{\#grantee}, \ldots &
  \code{acl:Authorization} (modes \code{acl:Read}, \code{acl:Write},
  \code{acl:Control}; Figure~\ref{fig:acls}) \\
\end{tabular}
\caption{Principal RDF classes per document kind. Prefixes are as registered
at \url{https://prefix.cc}, except: \code{rec:} is RealEstateCore
(\code{https://w3id.org/rec\#}), \code{geo:} is WGS84 lat/long
(\code{http://www.w3.org/2003/01/geo/wgs84\_pos\#}), and \code{bldg:},
\code{cons:}, \code{gran:} are application-owned
(\code{https://solid.ti.rw.fau.de/gra/\{building,consumption,vocab\}.ttl\#}).}
\label{tab:classes}
\end{table}

\section{The write path: storage models and projections}\label{sec:writes}

\subsection{Two storage models}\label{sec:models}

Every mutation commits against one of two models, chosen by \textbf{where
authority (ground truth) lives}:

\begin{description}
\item[In-place resource (the default).] The resource \emph{is} the state:
  \code{GET} $\to$ mutate $\to$ conditional \code{PUT} with \code{If-Match} ---
  optimistic concurrency in its plain HTTP form, retried through a re-read on
  \code{412}/\code{409} and surfaced as a conflict after bounded attempts.
  The conditionality is \emph{validator-dependent}: \code{If-Match} is sent when
  the read supplied an \code{ETag}; on servers that omit validators
  (\S\ref{sec:findings}) the write proceeds unconditioned and the model
  degrades, knowingly, to last-writer-wins --- acceptable for the data this
  model is reserved for (single-writer personal state, where the residual race
  is the same user's own devices), but a degradation the conformance findings
  make explicit rather than the design hiding it. Correct wherever the data has
  a single writer who owns it: building documents, energy datasets, preference
  and bookmark files, aggregation definitions and snapshots, the profile
  document.
\item[Event-sourced log (the escalation).] Chosen when in-place breaks down:
  several agents write the same state (concurrent \code{PUT}s would clobber), or
  the history is itself the point (audit, replay). Immutable events are
  \code{POST}ed to an LDP container --- the server mints each child IRI, so
  concurrent appends never collide --- and are never updated. The log is ground
  truth; \emph{current state is always derived through a projection}
  (\S\ref{sec:projections}). Instances: the outgoing and incoming share logs,
  per-room membership/role logs, the notification inbox.
\end{description}

The governing rule is one sentence: \emph{overwrite single-writer owned state in
place; event-source anything cross-agent or needing replay; treat enforcement
artifacts as projections of the log.}

Within the event model there are two \emph{delivery topologies}, chosen by
whether the affected party can read where the fact lives. A \textbf{shared
container} (the rooms) is one log all participants can read and append to;
nobody is notified --- readers fold when they look. An \textbf{inbox push} (the
sharing events) is used when the authoritative record sits behind the writer's
access control: the grant lives in the sharer's log and ACLs, which the
recipient cannot read, so a \emph{copy} of the event is delivered into the
recipient's inbox --- transport, not a log; the drain archives each message into
the recipient's own incoming log and deletes it, leaving the event recorded once
per party.

\subsection{Idempotent creation under a retrying transport}\label{sec:idempotency}

Every write rides a retry layer (transient throttling is a fact of hosted pods,
\S\ref{sec:findings}), which forces a creation discipline. Addressable resources
are created by \code{PUT} to a \emph{client-chosen} URI, not by
\code{POST}-and-\code{Location}: a \code{PUT} retry hits the same URI, and
\code{If-None-Match:\,*} turns a duplicate replay into a clean \code{412} --- the
idempotent-create primitive HTTP already has. A \code{POST} retry, by contrast,
mints a second slug on a lost \code{201} and silently duplicates the resource.
\code{POST} is reserved for exactly the places where a fresh server-minted child
per append is wanted and a retry duplicate folds away harmlessly: the event
logs. ``Harmlessly'' is a requirement on the folds, not luck: every projection
folds latest-per-key (a grant by (resource, grantee), a membership by member),
so a duplicated event coincides with its original --- appends are effectively
idempotent \emph{modulo the fold}, which is the algebraic condition retry
safety actually needs.
Rule of thumb: resources you must \emph{address later} get client URIs; resources
you only \emph{accumulate} get \code{POST}. Client-chosen URIs also let one flow
derive every dependent path (a building's datasets, a room's ACL and first join
event) without a write--parse--derive round-trip, and behave identically across
server implementations where \code{Slug}/\code{Location} semantics vary.

\subsection{Failure atomicity: ordering, not transactions}\label{sec:atomicity}

There are no multi-resource transactions, so a mutation that spans resources
(a building document plus its datasets; a grant's log event, ACL projection,
and inbox notification) states its failure model through \emph{write
ordering}: dependent resources first, the \emph{discoverable} resource last,
so the last write is the commit point and a crash before it leaves only inert
orphans no read will find. Bulk operations extend this per item: seeding $N$
demo entities is per-item best-effort with a tally outcome
(``added $n$ of $N$'') --- partial success is a \emph{result} the caller
renders, not an error --- and a retry mints fresh identifiers rather than
resuming. There is no compensation; what replaces rollback is the
reconciliation family of \S\ref{sec:triggers}: a grant recorded but not (or
no longer correctly) enforced is repaired by replay, a projection that drifted
is repaired by the audit's repair twin, and the archive restore \emph{embeds}
its reconciliation (the archive carries the log but not the derived ACLs, so
the restore mutation runs the replay as part of the same intent). Sagas
without compensation, ordered so that every interrupted state is either
invisible or repairable.

\subsection{Projections, and the enforcement cache}\label{sec:projections}

A log's derived state reaches its consumer through one of two disciplines, chosen
by \emph{whether the consumer can fold the log itself}:

\begin{description}
\item[Fold-on-read (the default).] The projection is computed in memory each
  time a query reads the log --- group events by key, keep the latest, emit
  current state --- and never persisted. The consumer is the application.
\item[Materialised projection.] Persisted as a store resource because the
  consumer \emph{cannot} fold. The sole instance is load-bearing: the WAC
  \code{.acl} files. The server's enforcement engine reads only \code{.acl}, so
  the projection must be written where the enforcer looks --- but the \code{.acl}
  is \emph{not} ground truth. It is an enforcement cache, rebuilt from the
  outgoing share log by a replay operation; a grant or revocation writes the ACL
  at the call site (mechanically an in-place \code{PUT}), while the authoritative
  record is the event. Every share dimension (recipient, energy included, which
  years, which attachments) is recorded \emph{in the event} precisely so the log
  stays self-sufficient for replay --- after an archive restore, which recovers
  the log but not the ACLs, the same replay rebuilds enforcement
  (obligation~1 of \S\ref{sec:running}).
\end{description}

\paragraph{Intensional scope, extensional projection.}
The projection's characteristic hazard is that a grant's recorded scope is
\emph{intensional} (``this building, all energy years'') while the applied ACLs
are \emph{extensional} (the per-dataset resources that existed at apply time). A
mutation that \emph{creates} a resource inside an intensionally granted scope is
a projection-input change with no projection update. Each resource-creating
mutation is audited against this hazard: attachment uploads are covered by
construction (the grant targets the container with a default rule, so later
members inherit); the energy-year write was the genuine gap and now runs a
best-effort re-application of the affected building's active grants per their
recorded scope (an all-years grant picks up the new dataset; a per-year grant
correctly leaves it outside); deletions are benign (a dropped resource takes its
ACL along, and replay skips grants on missing resources).

\begin{property}[Enforcement--intent agreement --- checkable, not invariant]\label{prop:agreement}
A recipient can read exactly what the folded outgoing log says they may.
\end{property}

This is a \emph{checkable agreement condition}, not an invariant: between an
in-scope resource creation and its (best-effort) reconciliation, and after a
silently failed repair, it can be violated, and no liveness is guaranteed ---
the system guarantees only that the audit \emph{detects} a violation and the
replay \emph{restores} agreement when run. The condition is executable: a headless test folds the log and probes each
implied resource with recipient-credentialed \code{GET}s; in the field the same
enumeration backs an audit operation (a dry-run diff of expected against actual
ACLs) and its repair twin (the replay). Auditability of the enforcement cache is,
we found, the single strongest argument for event-sourcing the grants: an
in-place permission model has nothing to diff against.

\paragraph{Read authority differs by direction.} Outgoing ``shared with whom''
reads the ACLs directly (a user can read their own) with the log as history;
incoming ``shared with me'' has no cheaper authority than the user's own incoming
log, folded on read --- and a missed revocation self-heals because the revoked
source \code{403}s on load and is pruned (\S\ref{sec:consistency}).

\subsection{The trigger axis, and the honest seams}\label{sec:triggers}

Orthogonally to mechanism, mutations split by \emph{trigger}: almost all are
\textbf{user-intent} (they record a decision a person made); a small family is
\textbf{reconciliation} --- system-initiated writes whose whole job is to make a
projection match observed reality: the stale-grant prune (appending a
self-revocation when a shared source has become inaccessible), the ACL replay,
and the grant re-application above. Reconciliations legitimately live inside
query and restore paths. That placement produces the system's three deliberate
violations of command--query separation~\cite{meyer} --- reads that may write:
the source load that prunes revoked shares, the inbox drain (a destructive move
shaped like a refresh), and a detail read that auto-materialises a missing
computed snapshot. Each is exceptional rather than per-call (the happy path
writes nothing), best-effort (a failed reconciliation degrades, never throws),
documented as a seam --- and, crucially for \S\ref{sec:mechanisms},
\emph{each still owns its cache effect}: the prune invalidates the
incoming-log key it changed, the user-invoked drain is an ordinary mutation
hook with its invalidations, and the snapshot auto-materialisation returns the
computed value as the read's own result, needing none. We prefer three named exceptions over either fiction:
pretending the reads are pure, or splitting each into a read plus a manual
cleanup step no user would run.

\section{The read path: a resource-oriented cache}\label{sec:reads}

\subsection{Mirror the graph, not the screens}\label{sec:normalised}

The client cache is \emph{resource-oriented}: one entry per store resource --- a
container listing, a building document, an energy dataset, an immutable log
event --- keyed by its IRI and namespaced by the viewer's WebID (a re-login
cannot read another user's cache). The domain shapes a screen wants (the building
list, the folded grant set, a labelled energy table) are \emph{never stored}:
they are derived at the edge, per render, by pure memoised selectors over the
resource entries. Four principles:

\begin{description}
\item[resource-keyed] the entry's value is parsed close to the triples, under
  canonical vocabulary keys --- no display labels in the cache.
\item[containers-as-queries] an LDP container is an entry over its membership
  listing; a \emph{structural} change (a resource added or removed) refreshes the
  container entry, a \emph{content} change refreshes the member's entry.
\item[derive-at-edge] selectors compose entries into app shapes in memory;
  presentation (labels, grouping, hiding) is applied at render, so a
  presentation-only change refetches nothing.
\item[single-read-path] every consumer --- map, dialog, background computation
  --- reads the \emph{same} entry; there is no second read to race or diverge
  from the warm copy.
\end{description}

This replaced a screen-shaped design in which one monolithic loader discovered,
fetched, parsed and composed many resources into a single view-shaped value under
a coarse key. Its failure mode was structural: a consumer needing a different
shape or fresher data re-read the store independently, and the independent read
raced the warm copy. Reads flow through four bands --- consumers; hooks with
combine selectors (which preserve the external
\code{\{data, isLoading, error\}} contract across internal reshaping); the
IRI-keyed resource entries; and React-free loader/parser functions shared with
the headless test path, so the application and its offline harness cannot drift
in how a resource is read.

\subsection{Structural invalidation}\label{sec:coverage}

The payoff of resource keying is that cache coverage stops being a per-call-site
discipline and becomes a property of the keying:

\begin{property}[Key coverage]\label{prop:coverage}
Let $K(r)$ be the set of resource entries a derived read $r$ depends on and
$W(m)$ the resources a write $m$ touches. Assume (i) $K(r) \supseteq
\mathit{inputs}(r)$ --- enforced structurally by derive-at-edge, since a
selector can only fold entries it subscribes to --- and (ii) every write path
of this client invalidates $W(m)$ soundly --- witnessed by the single key
registry that hooks and service peeks share. Then no write \emph{by this
client} changes $r$'s inputs without invalidating an entry in $K(r)$ or
changing $K(r)$ itself. (Writes by other agents are outside the property's
scope by construction; they are Guarantee~\ref{g:eventual}'s concern.)
\end{property}

The hazard this closes is concrete: the screen-shaped design once keyed a fold by
the \emph{identity of a set} (which buildings exist) but not the \emph{content it
folded} (which datasets each links), so a write that added an energy year to an
existing building changed no key and no refetch ever fired --- silently, because
last-good-wins rendering kept stale data on screen. Under resource keying the
write invalidates the touched dataset entry \emph{by IRI}; there is no whole-set
key left to under-cover. Three cost consequences are equally structural: a
content edit refetches one resource, not a screen; a structural change refetches
the container listing while member entries stay warm (the refold reads only the
new member); and a derived-state-only change (hiding a building) re-runs a
selector and refetches nothing. The principle: prefer correctness that
\emph{falls out of the data} --- keys derived from the folded resources --- over
correctness by discipline, every mutation remembering every dependent key, which
works until one call site forgets.

The trade-offs are the classic ones of normalised caching, accepted rather than
fought: more and smaller entries (structural sharing and garbage collection keep
them cheap; net fetch volume \emph{drops} because each resource is read once and
reused across consumers), and container-to-member waterfalls (bounded-concurrency
fan-out from the listing, not lazy per-render reads). One point stays
deliberately coarse: an archive restore may have replaced anything, so it
invalidates the entire cache and re-reads through the dependent chain, with
last-good-wins holding the pre-restore view until the maximal re-read lands.

\section{Consistency without push}\label{sec:consistency}

\subsection{Four mechanisms}\label{sec:mechanisms}

There is no push from the store, so freshness rests on four deliberate
mechanisms, centralised rather than per-screen:

\begin{enumerate}
\item \textbf{Server-driven freshness.} Every subscribed read revalidates with a
  conditional \code{GET} (zero client-side freshness window; an unchanged
  resource is a cheap \code{304}). The client never fabricates a TTL; the
  server's validator decides. Immutable resources --- log events, whose IRIs are
  server-minted and whose content never changes --- opt out entirely and are
  refreshed only by invalidation.
\item \textbf{Write-owned invalidation.} Every write path of this client ---
  a user-intent mutation \emph{or} a reconciliation (\S\ref{sec:triggers})
  --- \emph{owns} the keys it invalidates, declared against a single leaf
  registry of key builders that the service-side cache peeks share, so hook
  keys and service reads cannot drift. This includes the CQS seams: the prune
  invalidates the incoming-log key it appends to, the user-invoked drain is an
  ordinary mutation hook, and the snapshot auto-materialisation returns the
  computed value as the read's own result, needing none. No caller wires
  invalidation; forgetting is not an available error.
\item \textbf{Last-good-wins rendering.} An in-flight or failed refetch keeps
  the previous good projection on screen; staleness is never rendered as
  emptiness. (This is what makes silent under-coverage dangerous ---
  \S\ref{sec:coverage} --- and why coverage had to become structural.)
\item \textbf{Attention-triggered revalidation.} State appended by \emph{other}
  agents (a room log, an incoming share) can never be locally invalidated ---
  the writer is elsewhere. The client revalidates it when the user's attention
  lands on it: opening the room page invalidates the room log; visiting the
  sharing surface (or logging in) drains the inbox; a dialog's opening is a
  ``look'' that refetches what it displays. The \emph{mechanism} is standard
  client-side practice (data-fetching libraries ship refetch-on-focus and
  refetch-on-mount~\cite{reactquery}); the contribution is the policy and its
  contract --- attention as the \emph{only} poll trigger for cross-agent
  state, with the staleness bound of G\ref{g:eventual} --- frequent exactly
  where the user is, silent where they are not.
\end{enumerate}

Read-time reconciliation supplements these where a server's own reads lag its
writes: a delete polls the container listing until the member is gone; a load
that \code{403}s on a shared source prunes the corresponding grant
(\S\ref{sec:triggers}). The status taxonomy matters here: \code{403} is a
definitive access decision and \code{404} a definitive absence --- both prune
--- whereas a transport failure carries no status and must never prune;
treating an outage as a revocation would destroy records self-healing cannot
recover.

\subsection{The contract}\label{sec:contract}

To state the guarantees precisely we fix a minimal execution model. A
\emph{session} is one client's sequence of operation events: reads $r$ with
input set $\mathit{inputs}(r)$ (the resources folded) and a \emph{settle}
event (the response rendered), and writes $m$ with touch set $W(m)$ and a
settle event (the mutation reported done). The store assigns each resource a
per-resource version order; a read \emph{reflects} $m$ if, for every resource
in $\mathit{inputs}(r) \cap W(m)$, the version it serves is $\ge$ the version
$m$ installed. The guarantees are per session, in the style of session
guarantees~\cite{terry1994}, and they are \emph{safety} statements ---
\S\ref{sec:liveness} owns the liveness caveat.

\begin{guarantee}[Read-your-writes, own state]\label{g:ryw}
For a write $m$ and a read $r$ of the same session with
$\mathit{inputs}(r) \cap W(m) \neq \varnothing$: if $r$ \emph{settles} after
$m$ settles, $r$ reflects $m$. Mechanism: $m$'s owned invalidation marks the
touched entries before $m$ settles, so a later-settling $r$ cannot be served
from a pre-write entry; because the refetch is asynchronous, a render
\emph{between} $m$'s settle and $r$'s settle may still show the prior value,
marked as revalidating --- the guarantee is about settled reads, not
intermediate frames.
\end{guarantee}

\begin{guarantee}[Monotonic display]\label{g:monotonic}
Per cache entry, the version rendered is non-decreasing in the resource's
version order within a session: refetches replace state only on success
(last-good-wins), and immutable event entries never change. A cross-agent
revocation \emph{removes} an entry (a visibility change under
G\ref{g:heal}); removal is not a version regression.
\end{guarantee}

\begin{guarantee}[Attention-bounded eventual consistency, cross-agent]\label{g:eventual}
State written by another agent becomes visible no later than the next attention
event that covers it --- the page open, the inbox drain, the login refresh ---
and its staleness window is therefore bounded by the user's own interaction,
not by a timer.
\end{guarantee}

\begin{guarantee}[Self-healing visibility]\label{g:heal}
Access lost without notification (a revocation, a deleted source) is detected at
the next read of the affected resource (\code{403}/\code{404}), rendered
immediately, and reconciled into the local record by the prune --- the client
never depends on being told.
\end{guarantee}

Equally important is what is \emph{not} promised: no writes-follow-reads or
causal ordering across agents (two users editing unrelated resources have no
ordering at all; the event logs serialise only appends to one container); no
\emph{cross-entry snapshot} --- after a multi-resource mutation the touched
entries refetch independently, so a fold rendered mid-refetch can briefly mix
old and new entries (fractured reads; the settled state of G\ref{g:ryw} is
consistent, the intermediate frames are not); no bounded staleness in
wall-clock terms for a page the user never revisits; and no offline operation
(writes require the store; the cache is a projection, not a replica).

\paragraph{Safety and liveness.}\label{sec:liveness}
All four guarantees are safety-flavoured. Liveness is deliberately
conditional: G\ref{g:eventual}'s bound is an attention event that may never
occur, Property~\ref{prop:agreement}'s repair runs best-effort, and
reconciliations may fail silently and retry only on the next trigger. Making
that conditionality explicit is, we think, more honest than a wall-clock bound
no polling client can keep. We consider the narrowness of this contract a
feature: it is exactly what the mechanisms of \S\ref{sec:mechanisms} can
actually deliver, and each guarantee names the mechanism that delivers it.

\paragraph{Why not CRDTs.} The local-first program~\cite{kleppmann2019,
shapiro2011} solves a harder problem --- concurrent edits to shared replicas ---
that this storage discipline arranges not to have: in-place resources have a
single writer by design, and cross-agent state is append-only with commutative
folds (add/remove grant events keyed by subject; last-writer-wins per key at
fold time). Where a conflict could occur at all (two devices of the \emph{same}
user racing an in-place \code{PUT}), optimistic concurrency detects it and the
losing write re-reads --- rare enough, for personal master data, that detection
beats merge machinery. The trade is deliberate: no offline writes, in exchange
for a store whose ground truth is always a readable document rather than a
replica state.

\paragraph{Why not push.} The Solid Notifications
Protocol~\cite{solid-notifications} would replace G\ref{g:eventual}'s attention
bound with subscription latency, and we consider it the natural upgrade path.
It is not assumed because server support is uneven (R5), a browser client's
subscriptions die with its tab (wake-up still needs a look), and every
mechanism above remains necessary as the fallback on servers without it ---
push, when adopted, slots in as a fifth trigger for the same invalidation
machinery, not as a replacement architecture.

\subsection{Two caches, layered}\label{sec:twocaches}

The application cache is not the only cache in the loop: every request also
traverses the browser's HTTP cache, and the two differ in every dimension that
matters. The app cache is keyed by (viewer WebID, resource IRI) and holds
\emph{parsed} values; the browser cache is keyed by URL (plus \code{Vary}) and
holds \emph{representations}, with the server's validators deciding freshness.
The revalidating read deliberately composes them --- \code{cache: no-cache}
forces a conditional request while still letting a \code{304} be served from
the browser-cached body --- which is efficient exactly when the server's
validators are sound, and is how the content-negotiation unsoundness of
\S\ref{sec:findings} reached the application at all: the wrong body came from
the \emph{browser} cache, honestly revalidated against a broken validator. Two
consequences deserve naming. The layers must agree on identity: a cache keyed
by URL is oblivious to \emph{who} is asking, so on a shared device the browser
cache can retain representations across logins that the WebID-namespaced app
cache correctly separates --- private data should be served
\code{Cache-Control: private} (ideally \code{no-store} on shared machines), a
server-side obligation the client cannot substitute for. And the chattiness of
revalidate-always is a deliberate trade: the alternatives --- trusting
\code{Cache-Control} heuristics, container validators as digests for their
members, \code{HEAD} probes --- all reduce requests by trusting servers that
\S\ref{sec:findings} shows to be unreliable precisely in this machinery; we
revisit the trade under push (\S\ref{sec:limits}).

\section{Ecosystem findings}\label{sec:findings}

The design leans on ordinary HTTP machinery --- validators, conditional
requests, status codes --- and part of the contribution is reporting what three
Solid server implementations --- Community Solid Server (v5--v7), Node Solid
Server (v5), and a JavaScript Solid server (0.0.x) from our test matrix; exact
versions will be pinned in the final conformance table --- actually provide. Method: the
application instruments its authenticated fetch once at login, recording every
request's method, target, status, and timing into an in-app inspection log; the
hermetic test lanes run the same data layer against local server instances, so
a conformance regression fails a test rather than a user. The single-read-path
rule (\S\ref{sec:normalised}) is what makes this harness a valid instrument:
loaders and parsers are shared between the reactive hooks and the headless
orchestrators, and the hooks themselves --- including the cache layer, its
invalidation fan-out, and entry reuse --- run under the test runtime against a
real local server, so the quantitative measurements are taken on the
application's own read path, headless, with the browser needed only where the
browser's HTTP cache is itself the object of measurement
(\S\ref{sec:twocaches}).

\begin{description}
\item[Missing validators.] \code{GET} responses on one production server
  carried no \code{ETag}, silently degrading every conditional read to a full
  transfer --- the revalidation \emph{protocol} still behaves correctly
  (\code{staleTime}~0 + conditional \code{GET} simply never sees a \code{304}),
  but the efficiency claim of \S\ref{sec:mechanisms} is server-dependent, and a
  client cannot assume cheap revalidation on today's ecosystem.
\item[Unsound validators under content negotiation.] One implementation derived
  a container's \code{ETag} independently of the representation's media type: a
  Turtle conditional \code{GET} could receive \code{304 Not Modified} against a
  cached JSON-LD body --- a validator unsoundness that surfaces only through a
  browser cache, and that we fixed upstream by making the validator
  format-aware. HTTP already forbids this --- distinct representations of a
  resource are required to carry distinct validators~\cite{rfc9110} --- so the
  finding is not a gap in the protocol but a MUST that RDF servers, which
  negotiate formats freely, are evidently prone to violate.
\item[Retries below or above proof-of-possession.] Hosted pods sit behind CDN
  throttling (intermittent \code{429}/\code{503}). A retry layer placed
  \emph{below} DPoP signing replays a consumed proof and fails; the retry must
  wrap the authenticated fetch so every attempt is freshly signed. The same
  layer motivates the idempotent-creation discipline of \S\ref{sec:idempotency}.
\item[Discovery cost.] Index-free discovery is $N{+}1$ reads (list a container,
  \code{GET} each member). At the realistic scale of personal data
  (tens-to-hundreds of resources per collection) this is acceptable ---
  HTTP/2 multiplexing takes most of the sting out of the request count, leaving
  server-side work as the real cost --- and the normalised cache amortises it
  across consumers; the pattern to avoid is re-discovery per consumer, which
  the single-read-path rule exists to prevent.
\end{description}

A quantitative evaluation --- request counts and \code{304} rates per screen
under the normalised cache versus the screen-shaped predecessor, and measured
staleness windows for the cross-agent cases --- is planned on the existing
instrumentation and is not yet included in this draft.

\section{Related work}\label{sec:related}

\emph{Session guarantees and weak consistency.} Our contract adopts the
per-session framing of Terry et al.~\cite{terry1994}; the attention-bounded
guarantee (G\ref{g:eventual}) can be read as a session guarantee whose
anti-entropy trigger is the user's own navigation. Classic client--server
coherence assumes server callbacks or leases; we operate where neither exists.

\emph{Local-first and CRDTs.} Kleppmann et al.~\cite{kleppmann2019} argue for
replica-owning clients with CRDT merge~\cite{shapiro2011}; \S\ref{sec:contract}
states why this system deliberately stays on the other side of that line
(single-writer resources, append-only logs, no offline writes).

\emph{Event sourcing and materialised views.} The log-as-truth discipline is
event sourcing~\cite{fowler-es,helland2016}; the \code{.acl} maintenance is
materialised-view maintenance~\cite{gupta1995} in an unusual regime: the view's
only consumer is another party's \emph{enforcement engine}, the view must be
written into the store because the consumer cannot fold, and the
intensional/extensional gap of \S\ref{sec:projections} is the incremental
view-maintenance problem in access-control clothing.

\emph{HTTP caching.} The mechanisms build on standard conditional requests and
validators~\cite{rfc9111}; last-good-wins rendering is the application-level
cousin of \code{stale-while-revalidate}~\cite{rfc5861}. Client data-fetching
frameworks~\cite{reactquery} supply the cache substrate we key, and the
normalised shape has a practitioner lineage --- ORM identity maps, GraphQL
normalised stores such as Relay~\cite{relay} --- that resolves entities by
application-assigned ids; keying by \emph{dereferenceable IRIs} lets the
store's own graph play that role. The contributions here are what to key the
cache by (the LDP resource graph), who owns invalidation (every write path),
and what contract results.

\emph{Solid.} Prior Solid application work centres on authentication, access
control, and interoperability~\cite{solid,wac,solid-notifications}; we add the
data-management layer such applications need and report where the server
ecosystem falls short of the HTTP machinery it nominally inherits.

\section{Limitations and future work}\label{sec:limits}

The quantitative evaluation of \S\ref{sec:findings} is outstanding. The
contract is per-session and single-device; multi-device read-your-writes would
need either push or a device-spanning session token. The three CQS seams are
documented but not eliminated; the inbox drain in particular could become a pure
read over an idempotent archive. Optimistic UI updates are used in exactly one
place (room switching) and a principled account of when optimism is safe under
this contract is open. The execution model of \S\ref{sec:contract} is small
enough to mechanise (the cache, the invalidation rule, and G\ref{g:ryw} in a
model checker), and doing so --- including the fractured-read window the
contract excludes --- is the natural hardening step. Finally, adopting server push where available --- as a
fifth invalidation trigger with the attention mechanisms as fallback --- is
designed but unimplemented. Two operational details are likewise untreated:
container pagination (collections here stay small enough to list whole) and
write-response economy (\code{Prefer: return=minimal}); both matter one order
of magnitude up.

\section{Conclusion}

Deleting the database tier does not delete the database problems. An application
over personal data stores inherits concurrency control, cache coherence, and
view maintenance, and must solve them client-side over plain HTTP against
servers of varying fidelity. We showed a design that solves them systematically:
authority-driven storage models with a replayable enforcement projection,
a resource-keyed cache whose coverage is structural rather than disciplined,
and a four-mechanism freshness policy whose modest, explicit contract ---
read-your-writes at home, attention-bounded eventual consistency abroad,
self-healing visibility under revocation --- is, we argue, the honest
consistency level a decentralised web application can promise today.

\bibliographystyle{plain}
\bibliography{paper-client-database}

\end{document}
